\begin{table}[htp]
\centering

    % Esta primera línea define las columnas de la tabla. Los posibles tipos de columna son:
    % c: texto centrado
    % l: texto alineado a la izquierda
    % r: texto centrado a la derecha
    % p: columna de ancho fijo
    % Las columnas tienen ancho dinámico según la anchura máxima de los elementos que contengan.

    % Las columnas l/r/c no parten el texto en filas diferentes si éste es demasiado largo. Para ello, puede utilizar el tipo de columna de ancho fijo "p".
    
    % Las barras verticales | se usan para definir los bordes verticales de la tabla. Pruebe a eliminar algunas y observe qué ocurre.
    \begin{tabular}{ | p{2cm} | p{2cm} | p{3cm} | p{3cm} | p{2cm} | }
        
        % A continuación van las filas de la tabla. En cada fila, las columnas se separan con el carácter &
        % Para terminar una fila se usa \\
        % Para incluir un borde horizontal entre filas se usa \hline

        % Cabecera con textos en negrita:
        \hline
        \textbf{Motor} & \textbf{2D/3D} & \textbf{APIs} & \textbf{Objetivo} & \textbf{Licencia}\\
        \hline
        
        % Cuerpo de la tabla:
        Godot & Ambos & OpenGL, Vulkan y DirectX & Videojuegos: partículas, animaciones, luces y sombras  & MIT\\
        \hline
        Blender & Ambos & CUDA, OptiX, HIP, oneAPI y Metal  & Fotorrealismo, basado en físicas  & GPL\\
        \hline
        GTK & 2D & OpenGL y Vulkan & Interfaces de usuario & GPL\\
        \hline
        Unity & Ambos & OpenGL, Vulkan, DirectX y Metal & Videojuegos y animación  & Privativa\\
        \hline
        
    \end{tabular} 
    
    \caption{Comparación de motores de renderización}
    \label{tbl:comparison-engine} 
\end{table}